\section{De Prompting a SFT: cómo hacer que las buenas trayectorias queden naturalmente al frente}

CoT decoding admite que la trayectoria correcta puede estar oculta en lo profundo de la distribución y luego la encuentra con candidatos adicionales. Un objetivo más directo es cambiar la distribución misma: lograr que las respuestas con pasos, capaces de conducir a la respuesta correcta, tengan mayor probabilidad desde la primera generación. Tanto Prompting como supervised finetuning (SFT) reordenan la distribución, pero el primero modifica temporalmente la distribución condicional mediante el contexto, mientras que el segundo forma una preferencia duradera mediante la actualización de parámetros.

\lecturefigure{slide-12.jpg}{La siguiente pregunta: ¿puede una thoughtful response convertirse naturalmente en top-1?}{12}

\subsection{Qué cambian Few-shot CoT y el zero-shot trigger}

Few-shot CoT coloca en el prompt problemas similares con sus soluciones paso a paso, y el modelo imita a partir de ellos el formato de salida y la estrategia de resolución; «Let's think step by step», en cambio, usa una frase disparadora genérica para inducir un proceso intermedio más largo. Ninguno de los dos modifica los parámetros; su ventaja es que se despliegan rápido y pueden ajustarse a cada tarea, y sus desventajas son el espacio que ocupan en el contexto, la sensibilidad a la elección de ejemplos y que el desempeño suele verse afectado por la redacción y el orden.

\lecturefigure{slide-13.jpg}{Few-shot Chain-of-Thought y zero-shot «Let's think step by step»}{13}

\lecturefigure{slide-14.jpg}{Ventajas y limitaciones de Prompting: simple y eficaz, pero depende de los ejemplos}{14}

\lecturefigure{slide-15.jpg}{Crítica en clase: pedir que primero se vean ejemplos o repetir «piensa paso a paso» es una interfaz poco natural}{15}

Al leer estas tres páginas conviene distinguir entre la \textbf{información de la tarea} y las \textbf{señales de control del razonamiento}. Los ejemplos le indican al modelo el formato de salida, pero también pueden filtrar vocabulario del dominio, la dificultad y la distribución de respuestas; la frase disparadora genérica reduce la ingeniería por tarea, aunque por lo general rinde menos que un few-shot de alta calidad. Si solo se informa que «el prompt A es mejor que el prompt B», sin reportar los ejemplos, el orden, el presupuesto de tokens y la configuración de decodificación, no es posible determinar si la mejora proviene de una transferencia real de estrategia o de condiciones de contexto más favorables.

\teachervoice{00:12:40--00:16:20, el docente dice que few-shot CoT es muy eficaz, pero que si quien pregunta ya cuenta con problemas similares y sus soluciones completas, ¿para qué preguntarle al modelo?; y añadir cada vez «piensa paso a paso» tampoco parece una capacidad natural. El docente enfatiza: un prompt exitoso es una interfaz útil, pero eso no significa que el problema de entrenamiento esté resuelto.}

\begin{knowledgebox}{Términos clave: Prompting, In-Context Learning y Finetuning}
\textbf{Prompting} consiste en agregar instrucciones, ejemplos o estructura a la entrada; \textbf{in-context learning} es la adaptación del modelo a una tarea a partir del contexto sin actualizar parámetros; \textbf{finetuning}, en cambio, actualiza los parámetros mediante gradientes. El estado de los dos primeros suele desaparecer al terminar la solicitud, mientras que el último cambia la distribución base del modelo en solicitudes futuras.
\end{knowledgebox}

\subsection{SFT: convertir trayectorias humanas en un objetivo de máxima verosimilitud}

SFT recopila problemas $x$ y soluciones humanas paso a paso $y^*=(r^*,a^*)$, y luego maximiza la verosimilitud de la secuencia de referencia. Para el conjunto de datos $\mathcal D$, la pérdida es
\[
\mathcal L_{\text{SFT}}(\theta)
=-\mathbb E_{(x,y^*)\sim\mathcal D}
\left[\sum_{t=1}^{|y^*|}\log p_\theta(y_t^*\mid x,y_{<t}^*)\right].
\]
Aquí $\theta$ son los parámetros por actualizar y $|y^*|$ es la longitud de la respuesta de referencia. Con teacher forcing, durante el entrenamiento siempre se le da al modelo el prefijo correcto $y_{<t}^*$; en el despliegue, el modelo debe consumir el prefijo que él mismo generó, de modo que una desviación en un solo paso puede llevar los estados posteriores fuera de la distribución de entrenamiento.

\lecturefigure{slide-16.jpg}{Supervised Finetuning: recopilar pasos humanos y maximizar su verosimilitud}{16}

\lecturefigure{slide-17.jpg}{La brecha entre entrenamiento y prueba en SFT: de tareas anotadas a problemas nuevos}{17}

La figura 17 entrena con problemas de última letra y de manzanas, y luego prueba con «¿cuántas r tiene strawberry?». Aunque todas las tareas puedan escribirse como pasos en lenguaje natural, la coincidencia del formato superficial no garantiza la transferencia del algoritmo. El modelo puede aprender a «repetir las formas de las palabras y las estructuras de las oraciones de los ejemplos de entrenamiento» sin aprender a contar carácter por carácter en una cadena arbitraria; por eso se necesitan particiones out-of-distribution según la longitud de la combinación, la forma de las palabras, el dominio y la dificultad, en lugar de dividir al azar muestras parecidas.

\begin{warningbox}{El supuesto oculto de teacher forcing}
SFT supone que la trayectoria de referencia es correcta y, además, es un camino que el modelo genera con facilidad; también supone que la next-token imitation local produce un comportamiento globalmente correcto. Si el estilo de los pasos de referencia difiere demasiado de la distribución del modelo, este puede seguir sin entrar de manera estable en esa trayectoria aunque la pérdida de entrenamiento disminuya.
\end{warningbox}

\subsection{Por qué «recolectar más datos y ampliar la escala» no necesariamente resuelve la generalización}

La ventaja de SFT es su generalidad: una vez construido un formato unificado de entrada y respuesta, el mismo flujo de entrenamiento puede abarcar tareas de matemáticas, código y explicación. Su desventaja es que la generalización a combinaciones nuevas y a dificultades nuevas es inestable, y ampliar las trayectorias humanas del mismo tipo suele producir rendimientos decrecientes. El «Don't scale blindly» de la diapositiva no se opone a scale; exige confirmar primero que el objetivo de aprendizaje y el proceso de generación de datos estén alineados con el objetivo de despliegue.

\lecturefigure{slide-18.jpg}{SFT: general pero con generalización limitada; ampliar a ciegas el mismo paradigma no necesariamente funciona}{18}

\teachervoice{00:16:20--00:19:55, el docente recuerda que en 2021 su equipo descubrió que SFT generalizaba mal en razonamiento, y da una de las advertencias prácticas más fuertes de esta clase: si el paradigm en sí está equivocado, seguir ampliando datos del mismo tipo y la escala del entrenamiento no lo corregirá por sí solo.}

Tres tipos de prueba permiten distinguir entre «no aprendió lo suficiente» y «el objetivo está desalineado»:
\begin{enumerate}
  \item \textbf{Ampliación dentro de la misma distribución}: si agregar muestras parecidas a las de entrenamiento sigue produciendo mejoras estables;
  \item \textbf{Extrapolación combinatoria}: si el modelo colapsa cuando aumentan la longitud, el rango numérico o la combinación de operaciones;
  \item \textbf{Sustitución de trayectorias}: si soluciones distintas pero igualmente correctas se penalizan por error.
\end{enumerate}
Si el modelo solo mejora en el primer tipo, se comporta más bien como un interpolador; si en el tercer tipo trata trayectorias alternativas válidas como errores, el propio objetivo de máxima verosimilitud está comprimiendo el espacio de soluciones.

\subsection{Resumen de la sección}

Prompting induce trayectorias de forma temporal mediante el contexto, y SFT reordena la distribución de forma permanente mediante referencias humanas. Ambos pueden hacer que el reasoning aparezca con más frecuencia, pero los dos dependen de que una persona especifique «cómo debe verse». Cuando el problema de despliegue admite varios caminos válidos, su dificultad rebasa la distribución anotada o la trayectoria correcta es difícil de escribir a mano, el siguiente paso debe ser optimizar directamente la calidad del resultado, en lugar de seguir imitando un único camino.
